\subsection{不变元}

令 \(A[P]^W\) 为由 W 不变元素组成的 A{[}P{]} 的子代数。对于 \(p \in P\)，记 W.p 为 p 在 W 下的轨道，并令 \(S(e^p) = \sum_{q \in W \cdot p} e^q\) 为 \(e^p\) 的 W-变换之和；这是一个 W 不变元素。若 \(p \in P \cap \overline{C}\)，则 \(w(p) \leqslant p\)（对所有 \(w \in W\)）（第 1 节，第 6 号，命题 18），且 \(e^p\) 是 \(S(e^p)\) 的唯一极大项。

令 \(x = \sum_{p} x_{p} e^{p} \in \mathrm{A}[\mathrm{P}]^{\mathbb{W}}\)；有 \(x_{w(p)} = x_{p}\)（对所有 \(p \in \mathrm{P}\) 和 \(w \in \mathrm{W}\)）。另一方面，\(\mathbb{W}\) 在 \(\mathbb{P}\) 中的每个轨道都恰好与 \(\mathbb{P} \cap \overline{\mathbb{C}}\) 相交于一点（第 1 节，第 5 号定理 2）。因此，

\[
x = \sum_ {P \cap \overline {{C}}} x _ {p} S (e ^ {p}).\tag{6}
\]

我们得到：

引理 3. \(S(e^p)\)（其中 \(p \in P \cap \overline{C}\)）构成 A-模 \(A[P]^W\) 的一组基。

更一般地：

命题 3. 对任意 \(p \in \mathrm{P} \cap \overline{\mathrm{C}}\)，令 \(x_p\) 为 \(\mathrm{A}[\mathrm{P}]^{\mathrm{W}}\) 中以 \(e^p\) 为唯一极大项的元素。则族 \((x_p)_{p \in \mathrm{P} \cap \overline{\mathrm{C}}}\) 是 A-模 \(\mathrm{A}[\mathrm{P}]^{\mathrm{W}}\) 的一组基。

先证明一个引理。

引理 4. 设 I 是满足以下条件的有序集：

(MIN) I 的每个非空子集都含有一个极小元。

设 \(\mathbf{E}\) 是一个 A-模，\((e_i)_{i \in I}\) 是 \(\mathbf{E}\) 的一组基，\((x_i)_{i \in I}\) 是 \(\mathbf{E}\) 中满足下式的元素族：

\[
x _ {i} = e _ {i} + \sum_ {j <   i} a _ {i j} e _ {j},
\]

对所有 \(i \in \mathbf{I}\)（其中 \(a_{ij} \in \mathbf{A}\)，且族 \((a_{ij})\) 的支集对每个 \(i\) 都有限）。则 \((x_i)_{i \in \mathbf{I}}\) 是 \(\mathbf{E}\) 的一组基。

对于任意子集 \(J\)，它是 \(I\) 的子集；令 \(E_J\) 为 \(E\) 的子模，其基为 \((e_i)_{i \in J}\)。令 \(S\) 为 \(J\) 是 \(I\) 的子集且具有以下性质的对象所成的集合：

\begin{enumerate}
\def\labelenumi{(\alph{enumi})}
\item
  若 \(i' \leqslant i\) 且 \(i \in \mathbf{J}\)，则 \(i' \in \mathbf{J}\)；
\item
  \((x_{i})_{i\in J}\) 是 \(\mathrm{E_J}\) 的一组基。
\end{enumerate}

显然，按包含关系排序的 S 是归纳的且非空。因此它有极大元 J。若 \(J \neq I\)，令 \(i_{0}\) 为 I - J 的一个极小元，并置 \(J' = J \cup \{i_{0}\}\)。于是每个 \(i \in I\)（满足 \(i < i_{0}\)）都属于 J：由此 \(J'\) 满足 (a)。另一方面，\(J'\) 也满足 (b)：事实上，

\[
e _ {i _ {0}} = x _ {i _ {0}} - \sum_ {j <   i _ {0}} a _ {i _ {0} j} e _ {j},
\]

由此可得 (b)。所以 \(J' \in S\)，矛盾。故 J = I，引理得证。

现在证明命题 3。对 \(\mathbf{I} = \mathbf{P} \cap \overline{\mathbf{C}}\) 应用引理 4。令 \(q \in \mathbf{I}\)，并令 \(\mathbf{I}_q\) 为 \(p \in \mathbf{I}\) 中满足 \(p \leqslant q\) 的元素所成的集合。若 \(p \in \mathbf{I}_q\)，则关系

\[
q - p \geqslant 0, \quad p \in \overline {{\mathrm{C}}}, \quad q \in \overline {{\mathrm{C}}}
\]

蕴含

\[
(q - p | p) \geqslant 0 \text { and } (q - p | q) \geqslant 0,
\]

从而

\[
(p | p) \leqslant (p | q) \leqslant (q | q).
\]

因此集合 \(I_{q}\) 有界。由于 I 是离散的，故 \(I_{q}\) 有限，并且显然 I 满足条件 (MIN)。另一方面，对于所有 \(p \in I\)，

\[
x _ {p} = e ^ {p} + \sum_ {q <   p} c _ {p q} e ^ {q}
\]

所以根据 (6)，

\[
x _ {p} = \mathrm{S} (e ^ {p}) + \sum_ {q <   p, q \in \mathbb {I}} c _ {p q} \mathrm{S} (e ^ {q}).
\]

命题现在由引理 3 和引理 4 得出。

定理 1. 设 \(\overline{\omega}_1, \ldots, \overline{\omega}_l\) 是与室 C 对应的基本权，并且对于 \(1 \leqslant i \leqslant l\)，令 \(x_i\) 为 A{[}P{]} \(^W\) 中以 \(e^{\overline{\omega}_i}\) 为唯一极大项的元素。令

\[
\varphi : \Lambda [ X _ {1}, \dots , X _ {I} ] \rightarrow A [ P ] ^ {W}
\]

从多项式代数 \(A[X_{1},\ldots,X_{l}]\) 到 \(A[P]^{W}\) 的同态将 \(X_{i}\) 映到 \(x_{i}\)；则映射 \(\varphi\) 是同构。

由引理 2，\(\varphi\) 作用于单项式 \(X_{1}^{n_{1}}\ldots X_{l}^{n_{l}}\) 所得的像是具有唯一极大项 \(e^{n_1\overline{\omega}_1 + \dots +n_l\overline{\omega}_l}\) 的元素。由于 \(P\cap \overline{C}\) 的每个元素都能唯一写成 \(n_1\overline{\omega}_1 + \dots +n_l\overline{\omega}_l\) 的形式，命题 3 表明，\(\varphi\) 作用于单项式 \(X_{1}^{n_{1}}\ldots X_{l}^{n_{l}}\) 所得的像构成 \(A[P]^W\) 的一组基，定理得证。

例。1）可取 \(x_{i} = \mathrm{S}(\mathfrak{e}^{\overline{\omega}_i})\)。

\begin{enumerate}
\def\labelenumi{\arabic{enumi})}
\setcounter{enumi}{1}
\tightlist
\item
  根据第 3 号备注 2），可取 \(x_{i} = \mathrm{J}(e^{\rho +\overline{\omega}_i}) / d\)（记号同第 3 号）。
\end{enumerate}

